\section*{Introduction}
\addcontentsline{toc}{section}{Introduction}

Finite-state devices are a fundamental part of Computer Science. They find
applications in compiler design, natural language processing and many other fields.
Different flavours of finite-state devices have been developed to address specific
practical needs. In \cite{HMU01} Hopcroft and Ullman give extensive introduction to
finite-state automata. In \cite{MS19} Mihov and Schulz present an overview of a
generalized model of finite-state devices defined over monoids. In \cite{M97} Mohri
discusses applications of finite-state transducers in natural language processing.

\iffalse
Finite-state devices are part of a well-studied theoretical model that finds
applications in compiler design, natural language processing and many other fields.
They are discussed extensively in the works of Hopcroft and Ullman \cite{HMU01},
Mihov and Schulz \cite{MS19}, and others.

Over the years, many flavours of finite-state devices have been developed. One of
the fundamental classes of finite-state devices is the non-deterministic
finite-state automaton introduced by Rabin and Scott \cite{RS59}. In their classical
paper, Rabin and Scott prove that the expressive power of non-deterministic and
deterministic finite-state automata is the same. Furthermore, they outline an
effective algorithm that constructs a deterministic finite-state automaton
recognizing the same language as the one recognized by a source non-deterministic
finite-state automaton. The downside of this algorithm is that in the worst case the
number of states in the resulting deterministic finite-state automaton is
exponential in the number of states in the source non-deterministic finite-state
automaton.

In this thesis, we present a new kind of finite-state devices inspired by the
concept of graph separators. We show that for such devices there exist algorithms
for determinization and minimization that are more efficient than those applicable
to the general case. We discuss ways to combine these findings with existing
finite-state algorithms to achieve better practical results.

The rest of this work is organized as follows. In Section \ref{sec_preliminaries}
we provide formal preliminaries. Section \ref{sec_separators_finite_state_automata}
introduces the notion of separated finite-state automata and provides efficient
algorithms for determinization and minimization of such devices. Section
\ref{sec_separators_finite_state_transducers} contains similar results for separated
finite-state transducers.
\fi
